% Please do not change %
\documentclass[11pt]{article} % font size
\usepackage[margin=1in]{geometry} % margins
\usepackage{amsmath,amsthm,amssymb} % for the  common "math symbols"
\usepackage{algorithm}
\usepackage{algpseudocode}
\usepackage[shortlabels]{enumitem} % for enumeration
\usepackage{booktabs} % if you need tables
\usepackage{listings}
\usepackage{color}
\usepackage{quiver}
\usepackage{ulem}
\usepackage{wasysym}
\DeclareMathOperator*{\argmax}{argmax}
\DeclareMathOperator*{\argmin}{argmin}

% \theoremstyle{definition}
\newtheorem{theorem}{Theorem}[section]
\newtheorem{corollary}[theorem]{Corollary}
\newtheorem{lemma}[theorem]{Lemma}
\newtheorem{proposition}[theorem]{Proposition}
\newtheorem{claim}[theorem]{Claim}
\newtheorem{conjecture}[theorem]{Conjecture}
\theoremstyle{definition}
\newtheorem{definition}[theorem]{Definition}
\theoremstyle{definition}
\newtheorem{fact}[theorem]{Fact}
\theoremstyle{definition}
\newtheorem{example}[theorem]{Example}
\newtheorem{note}[theorem]{Note}

\lstset{
  frame=none,
  xleftmargin=2pt,
  stepnumber=1,
  numbers=left,
  numbersep=5pt,
  numberstyle=\ttfamily\tiny\color[gray]{0.3},
  belowcaptionskip=\bigskipamount,
  captionpos=b,
  escapeinside={*'}{'*},
  language=haskell,
  tabsize=2,
  emphstyle={\bf},
  commentstyle=\it,
  stringstyle=\mdseries\rmfamily,
  showspaces=false,
  keywordstyle=\bfseries\rmfamily,
  columns=flexible,
  basicstyle=\small\sffamily,
  showstringspaces=false,
  morecomment=[l]\%,
}
% Feel free to include additional packages (if needed), but make sure it does not override the margin/font requirements.
%
\setlength{\parindent}{0pt} % no indent for new paragraphs.
\setlength{\parskip}{5pt} % distance between paragraphs, which will be used when you have a blank line in your LaTeX code.


%%%%%%%%%%%%%%%%%%%%%%%%%%%%%%%%%%%%%%%%%%%%%%%%%%%%%%%%%%%%%%%%%%%%%%%%%%%%%%
\title{Categories, Proofs, and Processes: \\Lecture 17 Adjunctions}
%
\author{Wira Azmoon Ahmad}
\date{27 Feb 2024}
%
\begin{document}
%
\maketitle

%You may use the usual \LaTeX\ environments to structure your answers:
%\begin{align}
%        \label{basegnn}
%        \mathbf{h}_u^{(t+1)} = \sigma \bigg( \mathbf{W}_\text{self}^{(t)} \mathbf{h}_u^{(t)}  + \mathbf{W}_\text{neigh}^{(t)} \sum_{v \in N(u)} \mathbf{h}_v^{(t)} + \mathbf{b}^{(t)} \bigg): 
%\end{align}


%\begin{table}[h]
%\caption{A simple table.}
%\centering
%\begin{tabular}[t]{lrrrr}
%\toprule
%Graphs & Can & Be  & Fun & Too \\ 
%\midrule 
%$G_1$ & $1$  & $2$ & $3$ & $4$ \\ 
%$G_2$ & $-1$  & $-2$ & $-3$ & $-4$ \\ 
%$G_3$ & $0.1$  & $0.2$ & $0.3$ & $0.4$ \\ 
%\bottomrule
%\end{tabular}
%\label{tab}
%\end{table}

% A functor $F:\mathcal{C} \rightarrow \mathcal{D}$ is
% \begin{itemize}
%     \item \textbf{faithful} if each map $F_{A,B}:\mathcal{C}(A,B) \rightarrow \mathcal{D}(F\ A, F\ B)$ is injective;
%     \item \textbf{full} if each map $F_{A,B}:\mathcal{C}(A,B) \rightarrow \mathcal{D}(F\ A, F\ B)$ is surjective
% \end{itemize}


\addtocounter{section}{7}
\section{Monoidal Categories}
\begin{definition}[Informal]
    Monoidal categories are categories with a "parallel" composition $\otimes$.
\end{definition}
\subsection{Motivating Examples}
\begin{example}[$\textbf{Set}, \times$]
\[f: A\rightarrow B, \quad g: C\rightarrow D\]
\[f\times g: A\times C \rightarrow B\times D\]
\[(f\times g)(a,c) = (f(a),g(c))\]
% https://q.uiver.app/#q=WzAsNixbMSwwLCJBXFx0aW1lcyBDIl0sWzAsMCwiQSJdLFsyLDAsIkMiXSxbMCwxLCJCIl0sWzIsMSwiRCJdLFsxLDEsIkJcXHRpbWVzIEQiXSxbMCwxLCJcXHBpXzEiLDJdLFswLDIsIlxccGlfMiJdLFsxLDMsImYiLDJdLFsyLDQsImciXSxbNSwzLCJcXHBpXzEnIl0sWzUsNCwiXFxwaV8yJyIsMl0sWzAsNSwiZlxcdGltZXMgZyIsMCx7InN0eWxlIjp7ImJvZHkiOnsibmFtZSI6ImRhc2hlZCJ9fX1dXQ==
\[\begin{tikzcd}
	A & {A\times C} & C \\
	B & {B\times D} & D
	\arrow["{\pi_1}"', from=1-2, to=1-1]
	\arrow["{\pi_2}", from=1-2, to=1-3]
	\arrow["f"', from=1-1, to=2-1]
	\arrow["g", from=1-3, to=2-3]
	\arrow["{\pi_1'}", from=2-2, to=2-1]
	\arrow["{\pi_2'}"', from=2-2, to=2-3]
	\arrow["{f\times g}", dashed, from=1-2, to=2-2]
\end{tikzcd}\]


Any $\mathcal{C}$ with products
\[f\times g:= \langle f\circ \pi_1, g\circ \pi_2\rangle\]

\begin{equation}
    \begin{aligned}
        (f'\times g')\circ (f\times g) &= (f'\times g') \circ \langle f \circ \pi_1, g\ circ \pi_2\rangle\\
        &= \langle f' \circ f \circ \pi_1, g' \circ g \circ \pi_2\rangle\\
        &= (f' \circ f)\times (g' \circ g) \quad \text{(interchange law)}
    \end{aligned}
\end{equation}
\end{example}

\begin{example}[$\textbf{Set}, +$]
$+$ is the disjoint union.
\[f: A\rightarrow B, \quad g: C\rightarrow D\]
\[f+ g: A+ C \rightarrow B+ D\]
\[(f+ g)(a,c) = \begin{cases}
    f(x) & \text{if } x \in A\\
    g(x) & \text{if } x \in C\\
\end{cases}
\]

In $\mathcal{C}$ with coproducts:
\[f+g = [\iota_1 \circ f, \iota_2 \circ g]\]
"Co-tupling of maps"
% https://q.uiver.app/#q=WzAsNixbMSwwLCJBKyBDIl0sWzAsMCwiQSJdLFsyLDAsIkMiXSxbMCwxLCJCIl0sWzIsMSwiRCJdLFsxLDEsIkIrIEQiXSxbMSwwLCJcXGlvdGFfMSJdLFsyLDAsIlxcaW90YV8yIiwyXSxbMSwzLCJmIiwyXSxbMiw0LCJnIl0sWzMsNSwiXFxpb3RhXzEnIiwyXSxbNCw1LCJcXGlvdGFfMiciXSxbMCw1LCJmKyBnIiwwLHsic3R5bGUiOnsiYm9keSI6eyJuYW1lIjoiZGFzaGVkIn19fV1d
\[\begin{tikzcd}
	A & {A+ C} & C \\
	B & {B+ D} & D
	\arrow["{\iota_1}", from=1-1, to=1-2]
	\arrow["{\iota_2}"', from=1-3, to=1-2]
	\arrow["f"', from=1-1, to=2-1]
	\arrow["g", from=1-3, to=2-3]
	\arrow["{\iota_1'}"', from=2-1, to=2-2]
	\arrow["{\iota_2'}", from=2-3, to=2-2]
	\arrow["{f+ g}", dashed, from=1-2, to=2-2]
\end{tikzcd}\]
\[(f'+g') \circ (f+g) = (f'\circ f) + (g' \circ g)\]

\end{example}

\subsection{Non-Examples}
\begin{example}[$\textbf{Rel}, \times$]
$\times$ is not a product or a coproduct since relations are non-separable.
Let
\[\mathbb{B} = \{0,1\}\]
\[S = \{*\mapsto (0,0), * \mapsto (1,1)\}\]
\[R = \{*\mapsto 0, * \mapsto 1\}\]
% https://q.uiver.app/#q=WzAsNCxbMSwxLCJcXHsqXFx9Il0sWzEsMCwiXFxtYXRoYmJ7Qn1cXHRpbWVzXFxtYXRoYmJ7Qn0iXSxbMCwwLCJcXG1hdGhiYntCfSJdLFsyLDAsIlxcbWF0aGJie0J9Il0sWzAsMSwiUyIsMl0sWzEsMiwiXFxwaV8xIiwyXSxbMCwyXSxbMCwzXSxbMSwzLCJcXHBpXzIiXV0=
\[\begin{tikzcd}
	{\mathbb{B}} & {\mathbb{B}\times\mathbb{B}} & {\mathbb{B}} \\
	& {\{*\}}
	\arrow["S"', from=2-2, to=1-2]
	\arrow["{\pi_1}"', from=1-2, to=1-1]
	\arrow[from=2-2, to=1-1]
	\arrow[from=2-2, to=1-3]
	\arrow["{\pi_2}", from=1-2, to=1-3]
\end{tikzcd}\]
\[S' = \begin{cases}
    * \mapsto (0,0)\\
    * \mapsto (0,1)\\
    * \mapsto (1,0)\\
    * \mapsto (1,1)
\end{cases}
\]
$S'\neq S$, no unique relation.
\end{example}
\begin{example}[$\textbf{Vect}_{\mathbb{C}}, \otimes$]
$\otimes$ is the tensor product.
\end{example}

\subsection{Definitions}
\begin{definition}[Def. 114]
    A monoidal category $(\mathcal{C},\otimes, I)$ is a category $\mathcal{C}$ 
    with a functor $\otimes:\mathcal{C} \times \mathcal{C} \rightarrow \mathcal{C}$
    called the \underline{monoidal product}, 
    an object $I\in \text{Ob}(\mathcal{C})$ called the unit
    and natural isomorphisms
    \[\alpha_{A,B,C}: (A\otimes B)\otimes C \xrightarrow[]{\sim} A\otimes (B\otimes C)\]
    \[\lambda_A : I \otimes A \xrightarrow[]{\sim}, \quad \rho_A : A\otimes I \xrightarrow[]{\sim} A \]
    satisfying 2 \underline{coherence conditions}.
\end{definition}
\begin{example}
    \[\{(a,b),c | a,b,c\} \cong \{a,(b,c) | a,b,c\}\]
\end{example}
\begin{definition}
    A monoidal category is \underline{strict} if $\alpha,\lambda,\rho$ are all identities.
\end{definition}
\begin{claim}
Any monoidal category is equivalent to a strict category.
\end{claim}

\subsection{Aside: Term rewrite systems}
\begin{definition}
    A signature 
    \[\Sigma = \{\otimes: 2, I: 0\}\]
    where $\otimes$ is a symbol, $2,0$ are arities, with arity $0\Rightarrow$ constant.
\end{definition}
\begin{definition}
A term $t\in T_{\Sigma} (V)$ for some set $V$ (variables) 
is a tree whose nodes are in $\Sigma$ and leaves are in $V$.
\end{definition}
\begin{example}
    \[V = \{x,y,z\}, \quad T_{\Sigma} V \ni \begin{cases}
        (x \otimes x) \otimes x\\
        I \otimes I \\
        (x\otimes y) \otimes z\\
        ...
    \end{cases}\]
\end{example}
\begin{definition}
    A rewrite rule is a pair of terms
    \[(x\otimes y) \otimes z \xrightarrow[]{R_1}  x \otimes (y \otimes z)\]
    \[I \otimes x \xrightarrow[]{R_2} x\]
    \[x \otimes I \xrightarrow[]{R_3} x\]
    Rewrite rules can be applied by
\begin{enumerate}
    \item substituting variables
    \item replacing a sub-term
\end{enumerate}
For 
\[x\mapsto a, y\mapsto b, z\mapsto I\]
\[(a\otimes b)\otimes I \xrightarrow[]{R_1'} a\otimes (b\otimes I)\]
\[((a \otimes b)\otimes I) \otimes d \xrightarrow[]{R_1' \otimes d} (a\otimes (b\otimes I))\otimes d\]
\end{definition}
\begin{definition}
    "Good" rewrite systems are:
\begin{itemize}
    \item terminating (all sequences $t\rightarrow ... \rightarrow t'$ are finite)
    \item confluent:
    \[ \forall t_1,t_2. t\rightrightarrows\rightrightarrows\rightrightarrows t_1,t_2
    \exists t_3. t_1,t_2\rightrightarrows\rightrightarrows\rightrightarrows t_3\]
\end{itemize}
Terminating + confluent $\implies$ unique normal forms.
\end{definition}
\begin{example}[Confluence]
% https://q.uiver.app/#q=WzAsMyxbMCwwLCIoeFxcb3RpbWVzIEkpXFxvdGltZXMgeSJdLFswLDEsInhcXG90aW1lcyB5Il0sWzEsMCwieFxcb3RpbWVzIChJXFxvdGltZXMgeSkiXSxbMCwxXSxbMCwyXSxbMiwxXV0=
\[\begin{tikzcd}
	{(x\otimes I)\otimes y} & {x\otimes (I\otimes y)} \\
	{x\otimes y}
	\arrow[from=1-1, to=2-1]
	\arrow[from=1-1, to=1-2]
	\arrow[from=1-2, to=2-1]
\end{tikzcd}\]
\end{example}

\begin{definition}
"Coherence = confluence + commuting diagrams"
% https://q.uiver.app/#q=WzAsOCxbMCwyLCJUIl0sWzEsMV0sWzIsMCwiVF8xIl0sWzMsMV0sWzQsMiwiVF8zIl0sWzEsM10sWzIsNCwiVF8yIl0sWzMsM10sWzAsMV0sWzEsMl0sWzIsM10sWzMsNF0sWzAsNV0sWzUsNl0sWzYsN10sWzcsNF1d
\[\begin{tikzcd}
	&& {T_1} \\
	& {} && {} \\
	T &&&& {T_3} \\
	& {} && {} \\
	&& {T_2}
	\arrow[from=3-1, to=2-2]
	\arrow[from=2-2, to=1-3]
	\arrow[from=1-3, to=2-4]
	\arrow[from=2-4, to=3-5]
	\arrow[from=3-1, to=4-2]
	\arrow[from=4-2, to=5-3]
	\arrow[from=5-3, to=4-4]
	\arrow[from=4-4, to=3-5]
\end{tikzcd}\]
commutes.
\end{definition}

\begin{definition}[MacLane's coherence conditions]
% https://q.uiver.app/#q=WzAsOSxbMSwwLCIoQVxcb3RpbWVzIEkpIFxcb3RpbWVzIEIiXSxbMSwxLCJBXFxvdGltZXMgKEkgXFxvdGltZXMgQikiXSxbMiwxLCJBXFxvdGltZXMgQiJdLFswLDQsIigoQVxcb3RpbWVzIEIpXFxvdGltZXMgQylcXG90aW1lcyBEIl0sWzEsNSwiKEFcXG90aW1lcyAoQlxcb3RpbWVzIEMpKVxcb3RpbWVzIEQiXSxbMSwzLCIoQVxcb3RpbWVzIEIpXFxvdGltZXMgKENcXG90aW1lcyBEKSJdLFszLDQsIkFcXG90aW1lcyAoQlxcb3RpbWVzIChDXFxvdGltZXMgRCkpIl0sWzIsNSwiQVxcb3RpbWVzICgoQlxcb3RpbWVzIEMpXFxvdGltZXMgRCkiXSxbNSw0XSxbMCwxLCJcXGFscGhhIiwyXSxbMCwyLCJcXHJobyBcXG90aW1lcyBCIl0sWzEsMiwiQVxcb3RpbWVzIFxcbGFtYmRhIiwyXSxbMyw0LCJcXGFscGhhXFxvdGltZXMgRCIsMl0sWzMsNSwiXFxhbHBoYSJdLFs1LDYsIlxcYWxwaGFfe0FcXG90aW1lcyBCLCBDLER9Il0sWzQsNywiXFxhbHBoYSIsMl0sWzcsNiwiQVxcb3RpbWVzIFxcYWxwaGEiLDJdXQ==
\[\begin{tikzcd}
	& {(A\otimes I) \otimes B} \\
	& {A\otimes (I \otimes B)} & {A\otimes B} \\
	\\
	& {(A\otimes B)\otimes (C\otimes D)} \\
	{((A\otimes B)\otimes C)\otimes D} &&& {A\otimes (B\otimes (C\otimes D))} && {} \\
	& {(A\otimes (B\otimes C))\otimes D} & {A\otimes ((B\otimes C)\otimes D)}
	\arrow["\alpha"', from=1-2, to=2-2]
	\arrow["{\rho \otimes B}", from=1-2, to=2-3]
	\arrow["{A\otimes \lambda}"', from=2-2, to=2-3]
	\arrow["{\alpha\otimes D}"', from=5-1, to=6-2]
	\arrow["\alpha", from=5-1, to=4-2]
	\arrow["{\alpha_{A\otimes B, C,D}}", from=4-2, to=5-4]
	\arrow["\alpha"', from=6-2, to=6-3]
	\arrow["{A\otimes \alpha}"', from=6-3, to=5-4]
\end{tikzcd}\]
$(\Delta)$ (top) and $(\pentagon)$ (bottom) coherence.
\end{definition}
\begin{theorem}[Coherence]
    All formal$^*$ diagrams of $\lambda,\rho,\alpha$ commute.
\end{theorem}


\end{document}
